% Part: first-order-logic
% Chapter: introduction
% Section: syntax

\documentclass[../../../include/open-logic-section]{subfiles}

\begin{document}

\olfileid{fol}{int}{syn}

\olsection{句法}

我们首先必须精确界定哪些符号串算作一阶逻辑的语句。这要留待稍后进行；眼下我们暂且通过举例来推进。一阶逻辑的基本构件——词汇——分为两部分。第一部分是我们用来表达特定内容或指称特定事物的符号。我们用常项符号指称事物，并用谓词符号对我们所指称的事物作出陈述。例如，我们可以用 $\Obj a$ 作为常项符号来指称单个事物，然后用语句~$\Atom{\Obj P}{\Obj a}$ 对它作出某种断言。如果你心目中为“$\Obj a$”和“$\Obj P$”赋予了特定含义，你就可以将 $\Atom{\Obj P}{\Obj a}$ 读作一个英语句子（而且在你初次学习形式逻辑时，很可能已经这样做了）。有了一阶逻辑的此类简单语句后，你就可以利用词汇的第二部分——逻辑符号（联结词和量词）——来构建更复杂的语句。因此，例如，我们可以构造像 $(\Atom{\Obj P}{\Obj a} \land \Atom{\Obj Q}{\Obj b})$ 或~$\lexists[\Obj x][\Atom{\Obj P}{\Obj x}]$ 这样的表达式。

为了给语义及推导系统规则提供精确定义，以便进行严格的元逻辑研究，我们首先必须精确地定义什么算作一阶逻辑的语句。其基本思想很容易理解：我们可以仅凭谓词符号和常项符号形成一些简单语句，例如~$\Atom{\Obj P}{\Obj a}$。然后再利用联结词和量词，由这些简单语句形成更复杂的语句。但是，允许我们形成更复杂语句的规则究竟是什么？这些规则必须得到明确，否则我们对“一阶逻辑语句”的定义就不够精确。这里有几点问题。第一，要让正确的符号串算作语句；第二，要以一种使我们能够对\emph{所有}语句给出数学证明的方式来实现这一点。最后，我们还必须给出关于语句的一些基本操作的精确定义，例如“将~$!A$ 中的每个~$\Obj x$ 替换为~$\Obj b$”。麻烦在于，一阶逻辑中的量词和变元使得如何进行这一操作并非显而易见。例如，$\lexists[\Obj
x][\Atom{\Obj P}{\Obj a}]$ 应该算作一个语句吗？$\lexists[\Obj x][\lexists[\Obj x][\Atom{\Obj P}{\Obj x}]]$ 又如何？“将 $(\Atom{\Obj
P}{\Obj x} \land \lexists[\Obj x][\Atom{\Obj P}{\Obj x}])$ 中的 $\Obj x$ 替换为~$\Obj b$”的结果该是什么？

\end{document}